\chapter{Лекция}
\label{ch:intro}

\section{M-th power carrier recovery}

PSK сигнал можно представить в виде:

$$S(t) = Ae^{i(\omega_ct+\phi_k)}$$,

где A - амплитуда, $\omega_c$ - несущая, $\phi_k$ - фаза. \\

Фазу можно представить как $\frac{2\pi k}{M}$, где M - порядок модуляции (сколько точек в созвездии). Если у нас 8-PSK, то мы 
по сути на окружности расставляем 8 точек с равным интервалом $\frac{\pi}{4}$. \\

Подставим в сигнал фазу:

$$S(t) = Ae^{i(\omega_ct+\frac{2\pi k}{M})}$$

Предположим, что имеется частотный сдвиг $\triangle \omega_c$ между генератором на передатчике и генератором на приеме и фазовый сдвиг $\Theta$, тогда сигнал на приеме можно записать как:

$$S(t) = Ae^{i(\triangle \omega t+ \Theta + \frac{2\pi k}{M})}$$

Теперь делаем следующий фокус: возводим сигнал в степень M:

$$S(t)^M = A^Me^{M(i(\triangle \omega t+ \Theta + \frac{2\pi k}{M}))}$$

$$S(t)^M = A^Me^{(i(M\triangle \omega t+ M\Theta + M\frac{2\pi k}{M}))}$$

$$S(t)^M = A^Me^{(i(M\triangle \omega t+ M\Theta))} * e^{i(2\pi k)}$$

Заметим, что $e^{2\pi k}$ всегда равно 1, т.к $2\pi k$ всегда попадает в точку 0 на окружности. Тогда

$$S(t)^M = A^Me^{(i(M\triangle \omega t+ M\Theta))}$$

Мы избавились от модуляции и остался только частотный сдвиг $M\triangle \omega$, но мы не знаем, на какой частоте этот сдвиг находится.
Для определения частотного сдвига используется FFT. В месте частотного сдвига будет пик. \\

Когда мы найдем сдвиг, то нам останется только "повернуть" сигнал, домножив его на $e^{-2\pi \triangle fn}$, где n - номер отсчета. \\

Частоту, найденную в FFT, перед подстановкой в экспоненту нужно разделить на M, т.к мы строим спектр $M\triangle \omega$, а также
разделить частоту на sample rate SDR, тогда мы найдем сдвиг, который приходится на семпл. \\

Алгоритм:

\begin{enumerate}
    \item Возводим сигнал в степень M.
    \item Строим FFT.
    \item Ищем пик в FFT и фиксируем его частоту.
    \item Делим полученную частоту на M и Fs.
    \item Производим восстановление несущей.
\end{enumerate}

Данный метод устранит только большой сдвиг и вряд ли остановит вращение созвездия на приеме. Задача этого алгоритма - подготовить
сигнал к точной частотной синхронизации (fine frequency sync). Алгоритм FFS нельзя использовать сразу потому, что он способен восстановить
только мелкие сдвиги частоты, а на больших сдвигах работает плохо. Таким образом, два типа алгоритмов частотной синхронизации работают вместе:
сначала CFO, потом FFS для более тонкой коррекции. 

\section{Costas loop}

Петля Костаса - пример алгоритма точной частотной синхронизации. Работает после символьной синхронизации. Он основывается на созведии, т.е алгоритму должна быть известна используемая модуляция.
Коррекция происходит путем "подтягивания" символа к ближайшей "истинной" точке. \\

Пусть есть QPSK созвездие, на котором изображены "истинные" точки и одна точка, которую нужно "дотянуть" к "истинной":

\begin{figure}[H]
    \centering
    \includegraphics[width=1.0\textwidth]{costas.png}
    \caption{Петля Костаса}
\end{figure}

Угол $\alpha$ показывает разницу, на которую нужно "провернуть" точку, чтобы она встала на свое место. Задача сводится к нахождению
угла $\alpha$. \\

Предположим, что сигнал нормированный до [-1, 1], тогда $z = e^{i\beta}$, $a = e^{i\phi}$. Угол $\alpha$ - разница между углами
$\beta$ и $\phi$, т.е 

$$\alpha = \beta - \phi = e^{i(\beta - \phi)} = z \times a^*$$

В то же время угол $\alpha$ это $arctan \frac{Im{z\times a^*}}{Re{z\times a^*}}$. Arctan - дорогостоющая операция, но его можно аппроксимировать 
синусом при маленьких значениях. На единичной окружности синус любого угла это просто Im часть, потому что гипотенуза к любой точке 
на единичной окружности всегда равна 1, остается только противолежащий катет, который равен Im. В итоге получаем:

$$arctan \frac{Im{z}}{Re{z}} = sin(\alpha) = \sin(z\times a^*) = Im(z\times a^*) = \boxed{y\times a + b \times x}$$

Итоговое выражение сверху - ошибка. Далее ошибка подается на loop filter, где сглаживается и интегрируется.

\endinput
